\section{Related Work}\label{sec:related}

\paragraph{Lattice Boltzmann Methods for Multiphase Flows.}
The lattice Boltzmann method (LBM) has emerged as a powerful tool for simulating multiphase flows due to its simplicity in handling complex boundary conditions and natural parallelization. Several multiphase models exist: the pseudopotential model \citep{shan1993lattice}, the color-gradient model \citep{gunstensen1991lattice}, the free-energy model \citep{swift1995lattice}, and the phase-field model \citep{lee2005lattice}. Among these, phase-field approaches based on the Allen-Cahn equation \citep{fakhari2017improved} have gained popularity for their ability to handle high density ratios with good stability.

\paragraph{High Density Ratio LBM.}
Achieving stable simulations at high density ratios (e.g., water/air at 828:1) remains challenging. Early index-function models were limited to density ratios below 10. \citet{inamuro2004lattice} achieved density ratios of 1000 using a free-energy approach with Poisson pressure correction. \citet{lee2005pressure} developed a pressure-evolution model reaching density ratios of 1000. More recently, \citet{fakhari2017improved} proposed an improved locality-preserving AC-LBM that supports density ratios up to 1000 with good stability. Our work builds on this formulation with additional stabilization for curved boundaries.

\paragraph{Droplet Impact on Curved Surfaces.}
The study of droplet dynamics on curved substrates has attracted significant attention. \citet{liu2015symmetry} experimentally demonstrated asymmetric bouncing on \emph{Echeveria} leaves, showing 40\% contact time reduction due to momentum asymmetry. Their lattice Boltzmann simulations (using the Lee-Liu free-energy model) confirmed the mechanism but were limited to a single Weber number. \citet{xu2022droplet} studied droplet impact on triangular ridges using pseudopotential LBM, finding that contact time reduction depends on the interplay between Weber number and ridge geometry. \citet{xiao2025numerical} investigated droplet impact and freezing on supercooled curved surfaces using thermal LBM with a pseudopotential model, providing orthogonal experimental design analysis of We, wall temperature, contact angle, and curvature ratio.

\paragraph{Wetting Boundary Conditions on Curved Surfaces.}
Accurate wetting boundary conditions are critical for simulating droplet-surface interactions. \citet{fakhari2017diffuse} proposed a geometric wetting approach for curved boundaries in phase-field LBM. \citet{connington2015lattice} developed a wetting model for curved surfaces using the free-energy approach. Recent work by \citet{sashko2024phase} extended phase-field LBM wetting to 3D curved geometries with interpolation-based mitigation of staircase approximations. \citet{compact2025stencil} proposed a compact-stencil wetting BC for 3D curved surfaces, eliminating nonphysical droplet sliding and detachment.

\paragraph{Over-relaxation Techniques in LBM.}
Over-relaxation of boundary gradients is a known technique in LBM. \citet{ding2007lattice} used extrapolation-based wetting BCs with relaxation factors for free-energy LBM. \citet{zheng2020lattice} employed gradient interpolation with adjustable relaxation for phase-field models. However, these works focus on flat or mildly curved surfaces and do not address the specific resistance caused by the Allen-Cahn sharpening term on strongly curved hydrophobic surfaces. Our work differs in three ways: (1) we identify the AC sharpening as the specific source of resistance, (2) we derive $\alpha_{geo}$ from the AC-wetting balance (Eq.~\ref{eq:alpha_derivation}), and (3) we show that the optimal $\alpha_{geo}$ depends on $\theta_{eq}$, $R^*$, and $D_0/\xi$---a systematic calibration not previously available.

\paragraph{Gap.}
Despite progress, no existing study provides: (a) a derivation of the over-relaxation factor from the AC sharpening dynamics, (b) a systematic calibration across contact angle, curvature, and resolution, or (c) a comprehensive We $\times$ $R^*$ phase diagram for droplet impact on curved surfaces. Our work addresses all three.
